% =====================================================================
%  Conferenze, Zagarolo (On The Road), giorno 1 — Intervista 4
%  Formare game designer in una scuola di cinema: la NFTS
%  Fonte: trascrizione automatica "Track 12". Intervista in inglese ad
%  Andre, responsabile del corso di Games Design and Development della
%  National Film and Television School, condotta da un ex studente della
%  scuola, con molte domande dal pubblico (alcune in inglese incerto).
%  Etichette con prefisso "c2-4:". Nessun pacchetto aggiuntivo.
% =====================================================================

\chapter[Formare game designer: la NFTS]{Formare game designer in una scuola di
cinema: il modello della NFTS}
\label{cap:c2-4-nfts}

La \textbf{National Film and Television School} (NFTS) britannica è comparsa più
volte nei capitoli precedenti: vi hanno studiato l'ospite del
capitolo~\ref{cap:c2-1-marketing-publisher}, il producer del
capitolo~\ref{cap:c2-2-producer} vi ha lavorato per sei anni e mezzo, e lo stesso
moderatore di questa intervista è un ex studente. Molti relatori della conferenza
vengono da lì (il moderatore cita Manos, John e Tony, che parlano in altri
interventi). In questo incontro, in inglese, parla il responsabile del corso di
videogiochi, che spiega ``il segreto'' della scuola e risponde a molte domande
pratiche sull'ingresso nell'industria.

% ---------------------------------------------------------------------
\section{L'ospite}
\label{sec:c2-4-ospite}

\textbf{Andre}\footnote{Il cognome non è pronunciato.} è l'attuale
\textbf{responsabile del corso di Games Design and Development} della NFTS. Viene
dalle belle arti: ha studiato arte e si è insegnato da solo l'\textbf{animazione 3D}
su PC. Voleva dei colleghi, così si rivolse a un'agenzia per il lavoro in Olanda,
il suo paese d'origine, chiedendo chi usasse quelle competenze: gli indicarono
un'azienda di videogiochi. ``Non mi aspettavo di fare giochi, ma è successo'': ha
lavorato nei videogiochi per circa \textbf{quindici anni}, fino a diventare
\textbf{art director}, poi si è trasferito nel Regno Unito.

Interessato all'insegnamento, una delle prime istituzioni a cui scrisse fu proprio
la NFTS, che non aveva ancora un corso di videogiochi ma aveva annunciato di volerne
aprire uno: ``che cosa brillante, un corso di videogiochi in una scuola di cinema''.
Non ricevette risposta. Insegnò per circa dieci anni in una università
``tradizionale''; poi, dieci anni dopo, la NFTS lo contattò su LinkedIn per la
guida del corso. Ci lavora da tre anni.

% ---------------------------------------------------------------------
\section{Il ``segreto'' della NFTS}
\label{sec:c2-4-segreto}

Come fa la scuola a mandare tanti studenti nell'industria in tutto il mondo? Andre
indica due cose.

\subsection{È piccolissima}

La NFTS ha meno di mille studenti in tutto. Il corso di videogiochi dura
\textbf{due anni}, con \textbf{8--12 studenti per anno}. Ma ci sono moltissimi
docenti, quindi molto \textbf{supporto individuale}. Questo permette di accogliere
studenti con ambizioni e momenti di vita diversissimi:
\begin{itemize}
  \item chi ha già una laurea triennale in videogiochi e vuole approfondire un'area;
  \item chi è più grande, ha già fatto uno o due giochi e ``ha sbattuto contro un
        muro'' nel proprio sviluppo;
  \item chi viene da carriere completamente diverse: in una classe recente c'erano
        un'attrice, uno psicologo e una persona che aveva lavorato in un
        obitorio.
\end{itemize}
Mettere persone con prospettive così diverse in una stanza, attorno all'idea di
fare giochi, ``è una cosa molto potente''.

\subsection{È una scuola di cinema}

Il secondo motivo, quello per cui si era candidato: la scuola di cinema offre un
\textbf{ambiente creativo} in cui gli studenti di videogiochi collaborano con
specialisti di altri corsi: \textbf{sceneggiatori}, \textbf{compositori},
\textbf{sound designer}; e, tramite la fotogrammetria, con il reparto di
\textbf{modellismo}. Questo permette di fare qualcosa di \textbf{profondamente
personale}.

Nella sua università precedente non c'era un reparto di sound design: per la
colonna sonora gli studenti scaricavano musica gratuita da un sito. Alcune tracce
erano buone ma generiche, e Andre arrivò a scrivere nelle consegne che la musica
doveva venire \textbf{``oltre la prima pagina''} del sito, perché conosceva a memoria
le dieci tracce che comparivano sempre. Poter andare da uno studente di composizione
e avere una \textbf{colonna sonora su misura} è qualcosa di speciale.

% ---------------------------------------------------------------------
\section{La selezione: la spinta prima della tecnica}
\label{sec:c2-4-selezione}

Che cosa cerca in un candidato? Andre ``non si preoccupa troppo'' delle
\textbf{capacità tecniche}: è bello che un candidato abbia già fatto giochi, sappia
programmare o modellare in 3D, ma sono cose che la scuola sa insegnare bene. Cerca
invece \textbf{un desiderio reale}, e non generico (``voglio fare giochi''), ma
\textbf{specifico}. Chiede sempre:
\begin{quote}
  Che cosa pensi di poter dare al medium dei videogiochi? Che cosa vuoi esprimere,
  come persona?
\end{quote}
Quella spinta conta molto più della tecnica: ``se la spinta c'è, il resto
segue''.

% ---------------------------------------------------------------------
\section{Il rapporto con l'industria}
\label{sec:c2-4-industria}

Il moderatore racconta che quasi tutta la sua classe ha trovato lavoro subito. Andre
è più cauto: ``non rivendico quel livello di successo''. Ci sono anni ottimi e anni
più difficili, e \textbf{in questo momento è più difficile}, anche per i diplomati
NFTS, e non solo nei videogiochi ma in tutto il cinema, la televisione e l'industria
creativa.

Sul piano dei legami con l'industria:
\begin{itemize}
  \item Andre è l'\textbf{unico docente a tempo pieno}, insieme a un coordinatore
        che cura la parte organizzativa (``cruciale'');
  \item tutti gli altri docenti sono \textbf{part-time} e \textbf{lavorano
        nell'industria}: il tutoraggio individuale avviene dal punto di vista di
        chi fa giochi ogni giorno. Avere un problema tecnico e parlarne con un
        docente di programmazione che il giorno prima lavorava in un team su un
        gioco vero, e il giorno dopo ci tornerà, ``è ancora più importante che
        andare agli eventi'';
  \item comunque la scuola va a quanti più \textbf{eventi} possibile e incoraggia
        gli studenti a \textbf{mostrare} il proprio lavoro, perché ``il network è
        tutto''.
\end{itemize}

% ---------------------------------------------------------------------
\section{Le domande del pubblico}
\label{sec:c2-4-domande}

\subsection{Dall'informatica ai videogiochi}

Uno studente di informatica, interessato a motori di gioco e algoritmi più che ad
arte e design, nota che le convention sembrano concentrate sull'arte. Come
orientarsi? Andre: chi ha competenze o ambizioni di \textbf{programmazione a basso
livello} ``troverà lavoro'': è la strada più facile. ``Sono i designer e gli
artisti a faticare, in questo momento''; di programmatori ``non ce ne sono mai
abbastanza''. E soprattutto \textbf{parlare con le persone}: raccontare la propria
passione, condividere le competenze, connettersi.

\subsection{Che cosa farei se mi diplomassi quest'anno}

Una studentessa di design che si diploma quell'anno chiede che cosa farebbe Andre al
suo posto. Nel Regno Unito andrebbe a un \textbf{ritrovo mensile di sviluppatori} a
Londra\footnote{Trascritto ``Good Game Death Club''; il nome esatto non è
verificabile.} e comincerebbe a parlare con le persone. ``Trova qualcosa di simile
qui''\footnote{Segue nella trascrizione una lunga ripetizione meccanica senza
senso.}.

\subsection{Il portfolio}

Quali buone pratiche per un portfolio, e quali errori evitare? Andre, che da art
director i portfolio li valutava:
\begin{itemize}
  \item \textbf{evitare i soggetti abusati}: ``per favore, non aggiungete un altro
        Kalashnikov all'assortimento sempre crescente di
        Kalashnikov''\footnote{Trascritto ``Kolesnikovs''; più avanti, con lo stesso
        spirito, ``non un'auto che si schianta'' (\emph{collapsing card} nella
        trascrizione). Il senso è: evitare i soggetti più comuni nei portfolio 3D.},
        ``e non è una battuta''. Creare cose \textbf{originali}, uniche, è molto
        potente;
  \item renderlo \textbf{accessibile}: niente siti contorti con sottosezioni.
        ``Il tuo portfolio sarà uno dei quindici che devo guardare in quindici
        minuti il venerdì sera'': più è facile vedere che cosa conta per te e quali
        sono le tue competenze, meglio è. ``Semplice.''
\end{itemize}

\subsection{Farsi conoscere come musicista}

Uno studente di sound design e musica, che compone in generi come chiptune, 8 bit e
classica, poco richiesti rispetto al pop e all'hip hop dei giochi ``ad alta
energia'', chiede come costruirsi contatti fuori dall'università. La risposta: le
\textbf{game jam}, già citate più volte nella giornata. Sono un ottimo posto per
conoscere persone e \textbf{dimostrare le proprie competenze in un gioco vero}, che
per un audio designer è la vetrina migliore.

\subsection{Buttare dentro chi non sa nulla}

Una domanda in due parti. La prima: è problematico introdurre direttamente nello
sviluppo di videogiochi chi non sa nulla (per esempio un aspirante programmatore che
salta algoritmi e strutture dati e passa subito ai Blueprint di Unreal)?

``È esattamente quello che facciamo.'' Tutto, nel corso, ruota attorno al
\textbf{fare giochi}:
\begin{itemize}
  \item nel \textbf{primo anno} gli studenti fanno \textbf{cinque giochi}, nel
        \textbf{secondo uno}: alla fine hanno sei giochi in portfolio;
  \item il primo progetto si chiama \textbf{Hello World(s)}: due settimane di
        introduzione in cui modellano un barile 3D e fanno un semplice gioco in C\#
        con Unity;
  \item poi vengono ``buttati'' in un evento in cui devono presentare \textbf{tre
        concept} ai docenti, che scelgono il più fattibile;
  \item hanno \textbf{sei settimane} per realizzarlo, come team di 8--12 persone,
        ciascuno con una disciplina assegnata: project manager, lead programmer
        (``anche se non ha mai programmato in vita sua''), UI artist, narrative
        designer\dots ``E diciamo loro: arrangiatevi.''
\end{itemize}
È spaventoso per loro e anche per Andre, che non sa dove porterà; ma in tre anni,
tre volte, ne è uscito ``un'esperienza di gioco bellissima''. Funziona perché è tutto
\textbf{problem solving}: ``non so come si fa, va bene, scopriamolo''. In pedagogia
si chiama \textbf{apprendimento basato sui problemi} (\emph{problem-based learning}),
e Andre ci crede fermamente.

La seconda parte, sui ruoli meno noti dell'industria (marketing, produzione, QA),
resta senza risposta esplicita: Andre risponde alla prima e la domanda si
chiude\footnote{L'interlocutore gli dà facoltà di scegliere una delle due
domande.}.

\subsection{L'intelligenza artificiale}

Avete iniziato a usare l'IA nel flusso di lavoro? Andre trova ``incredibilmente
difficile'' avere un'opinione, perché ``ogni volta che penso qualcosa, due settimane
dopo è diverso''. La scuola ha dei \textbf{principi} sull'IA:
\begin{itemize}
  \item non è contraria all'uso dell'IA, \textbf{purché non tolga lavoro ad altre
        persone della scuola}. La collaborazione è il cuore della NFTS, tra il
        corso di videogiochi e tutti gli altri (regia, produzione televisiva\dots).
        Se l'IA sostituisce il compito di qualcuno, togliendo a uno studente
        un'opportunità di apprendimento, è un problema;
  \item in tutti gli altri casi, sperimentare con l'IA va bene, e i docenti
        condividono come la usano nel loro lavoro;
  \item il principio chiave: \textbf{le persone devono avere il controllo}. È uno
        strumento, ha un suo posto, ma serve qualcuno che lo governi.
\end{itemize}

\subsection{Specializzarsi o diversificarsi nel 3D}

Per entrare come modellatore 3D servono molti progetti o uno solidissimo? Andre:
\textbf{un progetto molto buono vale più di due medi}. Meglio tre oggetti 3D di
grande qualità (``non un'auto che si schianta, per favore''), mostrando non solo il
risultato finale ma anche \textbf{il processo}, comprese le revisioni: così si vede
come il candidato individua e risolve i problemi. Sulla \textbf{specializzazione}:
dipende dallo studio. Un grande studio tripla A cerca artisti
\textbf{specializzati} (``bravissimi nelle auto o nella vegetazione''); uno studio
piccolo probabilmente vorrà un portfolio \textbf{più vario}, perché si fanno molte
cose diverse.

\subsection{Il software preferito}

L'ultima domanda: quali programmi preferisce? Andre ha attraversato tutta la
storia: \textbf{3D Studio} (su DOS), poi \textbf{3D Studio Max}, poi \textbf{Maya}
quando il suo studio fu ingaggiato da Sony, che lavorava solo con Maya; ora sta
imparando \textbf{Blender} perché è ciò che la maggior parte degli studenti vuole
usare. ``È solo software.'' Contano le \textbf{competenze di base}, qualunque sia lo
strumento:
\begin{itemize}
  \item sapere che aspetto ha una \textbf{buona topologia};
  \item saper fare l'\textbf{UV unwrapping}, ``anche se l'IA lo fa per te: devi
        comunque entrare a correggere i problemi'';
  \item saper \textbf{risolvere problemi};
  \item saper \textbf{scomporre una forma complessa in primitive 3D}.
\end{itemize}
Allo stesso modo, conclude, un programmatore direbbe che l'IDE non fa una grande
differenza nel programma.

% ---------------------------------------------------------------------
\section*{Punti chiave del capitolo}
\addcontentsline{toc}{section}{Punti chiave del capitolo}

\begin{itemize}
  \item La forza della NFTS: classi \textbf{piccolissime} (8--12 per anno) con molto
        tutoraggio individuale, studenti dai percorsi più vari, e la
        \textbf{collaborazione} con sceneggiatori, compositori, sound designer e
        modellisti di una scuola di cinema.
  \item In selezione conta la \textbf{spinta specifica} (``che cosa vuoi dare al
        medium?'') più della tecnica, che si insegna.
  \item Docenti part-time \textbf{attivi nell'industria}; il momento attuale è
        difficile per tutti i diplomati delle industrie creative.
  \item I programmatori trovano lavoro più facilmente di designer e artisti; per
        tutti, \textbf{parlare con le persone}, frequentare i ritrovi di
        sviluppatori, fare \textbf{game jam}.
  \item Portfolio: originale (niente soggetti abusati), \textbf{semplice da
        consultare}, pochi lavori eccellenti con il \textbf{processo} in vista;
        specializzato per i grandi studi, vario per i piccoli.
  \item Didattica: si impara facendo giochi (sei in due anni), con ruoli assegnati
        anche a chi non li conosce: \textbf{problem-based learning}.
  \item IA: ammessa se non sottrae lavoro e opportunità di apprendimento agli
        altri; \textbf{le persone devono avere il controllo}.
  \item Il software cambia (3D Studio, Max, Maya, Blender): restano topologia, UV,
        problem solving, scomposizione delle forme.
\end{itemize}
